\documentclass[12pt, a4paper]{article}

% 引入必要的宏包
\usepackage[UTF8]{ctex}     % 中文支持
\usepackage{amsmath, amssymb, amsthm} % 数学符号和环境
\usepackage{geometry}       % 页面设置
\usepackage{enumitem}       % 列表定制
\usepackage{graphicx}       % 图片支持
\usepackage{fancyhdr}       % 页眉页脚

% 页面边距设置
\geometry{left=2.5cm, right=2.5cm, top=2.5cm, bottom=2.5cm}

% 宏定义：方便排版选择题选项
% 使用方法：\choices{选项A}{选项B}{选项C}{选项D}
\newcommand{\choices}[4]{
    \begin{enumerate}[label=\Alph*., itemsep=0pt, parsep=0pt, topsep=5pt, labelsep=0.5em]
        \item #1
        \item #2
        \item #3
        \item #4
    \end{enumerate}
}
% 横向排列的选项 (2列)
\newcommand{\choicesTwo}[4]{
    \begin{enumerate}[label=\Alph*., itemsep=0pt, parsep=0pt, topsep=5pt, labelsep=0.5em]
        \begin{minipage}{0.45\linewidth} \item #1 \end{minipage}
        \begin{minipage}{0.45\linewidth} \item #2 \end{minipage}
        \begin{minipage}{0.45\linewidth} \item #3 \end{minipage}
        \begin{minipage}{0.45\linewidth} \item #4 \end{minipage}
    \end{enumerate}
}
% 横向排列的选项 (4列)
\newcommand{\choicesFour}[4]{
    \begin{enumerate}[label=\Alph*., itemsep=0pt, parsep=0pt, topsep=5pt, labelsep=0.5em]
        \begin{minipage}{0.22\linewidth} \item #1 \end{minipage}
        \begin{minipage}{0.22\linewidth} \item #2 \end{minipage}
        \begin{minipage}{0.22\linewidth} \item #3 \end{minipage}
        \begin{minipage}{0.22\linewidth} \item #4 \end{minipage}
    \end{enumerate}
}

\title{\textbf{《高等数学》必刷习题——基础版}}
\author{}
\date{}

\begin{document}

\section*{第十章 向量代数与空间解析几何（基础）}

\begin{enumerate}
    \item 已知点 $ A(2,2,\sqrt{2}) $ , $ B(1,3,0) $ ,则向量 $ \overrightarrow{AB} $ 的模和方向角为 \underline{\hspace{1cm}}。
    \choicesTwo
    {$2, (\frac{2\pi}{3},\frac{\pi}{3},\frac{\pi}{4})$}
    {$1, (\frac{2\pi}{3},\frac{\pi}{3},\frac{3\pi}{4})$}
    {$2, (\frac{2\pi}{3},\frac{\pi}{3},\frac{3\pi}{4})$}
    {$1, (\frac{2\pi}{3},\frac{\pi}{3},\frac{\pi}{4})$}

    \item 设数 $ \lambda_{1}, \lambda_{2}, \lambda_{3} $ 不全为零，使得 $ \lambda_{1}\mathbf{a} + \lambda_{2}\mathbf{b} + \lambda_{3}\mathbf{c} = 0 $ ，则 $\mathbf{a}, \mathbf{b}, \mathbf{c}$ 三向量 \underline{\hspace{1cm}}。
    \choicesFour{线性无关}{共线}{共面}{共点}

    \item 下列各平面中，与平面 $ x + 2y - 3z = 6 $ 垂直的是 \underline{\hspace{1cm}}。
    \choicesTwo
    {$ 2x+4y-6z=1 $}
    {$ 2x + 4y - 6z = 12 $}
    {$ \frac{x}{-1}+\frac{y}{2}+\frac{z}{3}=1 $}
    {$ -x+2y+z=1 $}

    \item 直线 $\begin{cases}x+y+z-4=0\\ x-y-z+2=0\end{cases}$ 与直线 $\begin{cases}x-2y-z-1=0\\ x-y-2z=0\end{cases}$ 的位置关系 \underline{\hspace{1cm}}。

    \item 空间曲线 $ x=t,y=t^{2},z=t^{3} $ 在点 $(1,1,1)$ 的法平面方程 \underline{\hspace{2cm}}。

    \item 已知两点 $ A(-1,2,0) $ 和 $ B(2,-3,\sqrt{2}) $ ，则与向量 $ \overrightarrow{AB} $ 同方向的单位向量为 \underline{\hspace{2cm}}。

    \item 设 $ \mathbf{a} \neq 0 $ , 则与向量 $\mathbf{a}$ 同方向的单位向量 $ \mathbf{e} = $ \underline{\hspace{2cm}}。

    \item 设已知两点 $ M_{1}(2,2,\sqrt{2}) $ 和 $ M_{2}(1,3,0) $ ，则向量的模 $ \left|\overrightarrow{M_{1}M_{2}}\right|= $ \underline{\hspace{2cm}}。

    \item 向量 $ (1,0,1) $ 与 $ (1,\sqrt{2},1) $ 的夹角是 \underline{\hspace{2cm}}。

    \item 设 $ \mathbf{a} = \{1, 2, 3\} $ , $ \mathbf{b} = \{0, 1, -2\} $ , 则 $ (\mathbf{a} + \mathbf{b}) \times (\mathbf{a} - \mathbf{b}) = $ \underline{\hspace{2cm}}。

    \item 向量 $ \left|\mathbf{a}\right|=13,\left|\mathbf{b}\right|=19,\left|\mathbf{a}+\mathbf{b}\right|=24 $ , 则 $ \left|\mathbf{a}-\mathbf{b}\right|= $ \underline{\hspace{2cm}}。

    \item 点 $ (1,2,3) $ 到平面 $ 2x + y - 2z + 4 = 0 $ 的距离是 \underline{\hspace{2cm}}。

    \item 通过三点 $ (0,0,0) $ , $ (1,0,1) $ , $ (2,1,0) $ 的平面方程是 \underline{\hspace{2cm}}。

    \item 过点 $ (1,0,1) $ 和平面 $ x+y-z-1=0 $ 、平面 $ 2x+z+1=0 $ 平行的直线方程。

    \item 求通过点 $ M_{1}(3,-5,1) $ ， $ M_{2}(4,1,2) $ 且垂直于平面 $ x-8y+3z-1=0 $ 的平面方程。
\end{enumerate}


\end{document}